% Vietnamese translation for OI Public Library
% Source-PDF: IOI中国国家候选队论文/2005/胡伟栋/附件/Intuitionistic Logic.pdf
\documentclass[11pt]{article}
\usepackage{oipl-vn}

\title{Logic trực giác}
\author{ACM ICPC 2002--2003, Northeastern European Region, Northern Subregion}
\date{}

\begin{document}
\maketitle

\origtitle{Problem G. Intuitionistic Logic}
\sourcepath{IOI Candidate Team Papers/2005/Hu Weidong/appendix/Intuitionistic Logic.pdf}

\begin{center}
\begin{tabular}{rl}
Tệp vào: & \texttt{logic.in}\\
Tệp ra: & \texttt{logic.out}\\
Giới hạn thời gian: & 5 giây\\
Giới hạn bộ nhớ: & 64 megabyte
\end{tabular}
\end{center}

Gần đây Vasya làm quen với một trào lưu thú vị trong toán học và logic gọi là
``chủ nghĩa trực giác''. Ý tưởng chính của trào lưu này là bác bỏ luật bài
trung, tức luật logic nói rằng mọi mệnh đề đều hoặc đúng hoặc sai. Vasya rất
thích ý tưởng đó; cậu nói: ``Toán học cổ điển nói rằng định lý cuối cùng của
Fermat hoặc đúng hoặc sai, nhưng phát biểu này hoàn toàn vô dụng với tôi cho
đến khi tôi thấy một chứng minh hoặc một phản ví dụ''. Vì thế Vasya trở thành
một người theo logic trực giác.

Cậu cố dùng logic trực giác trong mọi hoạt động của mình, đặc biệt là trong
công việc khoa học. Nhưng logic này khó hơn logic cổ điển rất nhiều. Vasya
thường thử dùng những công thức logic đúng trong logic cổ điển nhưng không
đúng trong logic trực giác.

Bây giờ cậu muốn viết một chương trình giúp tự động kiểm tra tính đúng đắn của
các công thức. Cậu đã tìm được một cuốn sách mô tả cách làm, nhưng không giỏi
lập trình, nên bạn phải giúp cậu.

Việc xây dựng bắt đầu từ một đồ thị có hướng không chu trình tùy ý
\(\mathcal X=\langle X,G\rangle\). Sau đó ta xây dựng một thứ tự bộ phận trên
\(X\), tập các đỉnh của \(\mathcal X\): với mọi \(x,y\in X\), định nghĩa
\(x\le y\) khi và chỉ khi tồn tại một đường đi, có thể có độ dài \(0\), trong
\(\mathcal X\) từ \(x\) đến \(y\). Tiếp theo, xét tập \(P\) gồm mọi tập con của
\(X\), và tập \(H\subset P\) gồm mọi \(\alpha\subset X\) sao cho hai phần tử
khác nhau bất kỳ \(x,y\in\alpha\) đều không so sánh được, nghĩa là không có
\(x\le y\) và cũng không có \(y\le x\). Lưu ý rằng \(H\) luôn chứa tập rỗng và
mọi tập con một phần tử của \(X\).

Bây giờ ta có thể định nghĩa ánh xạ \(\operatorname{Max}:P\to H\subset P\).
Với mọi \(M\subset X\), đặt
\[
  \operatorname{Max}(M)
  =\{x\in M:\nexists y\in M,\ x\ne y,\ x\le y\},
\]
tức là tập tất cả các phần tử cực đại của \(M\).

Tiếp theo ta định nghĩa một số phép toán trên \(H\). Với mọi
\(\alpha,\beta\in H\), đặt
\[
  \alpha\Rightarrow\beta
  =\{x\in\beta:\nexists y\in\alpha,\ x\le y\},
\]
\[
  \alpha\land\beta=\operatorname{Max}(\alpha\cup\beta),
\]
\[
  \alpha\lor\beta
  =\operatorname{Max}(\{x\in X:\exists y\in\alpha,\ z\in\beta,\ x\le y,\ x\le z\}),
\]
\[
  0=\operatorname{Max}(X),\qquad 1=\varnothing,
\]
\[
  \neg\alpha=(\alpha\Rightarrow 0),\qquad
  \alpha\equiv\beta=((\alpha\Rightarrow\beta)\land(\beta\Rightarrow\alpha)).
\]

Bây giờ xét các công thức logic gồm những ký hiệu sau:

\begin{itemize}
  \item Hằng số \(1\) và \(0\).
  \item Biến: các chữ cái in hoa từ \texttt{A} đến \texttt{Z}.
  \item Dấu ngoặc: nếu \(E\) là một công thức, thì \((E)\) cũng là một công
  thức.
  \item Phủ định: \(\neg E\) là một công thức với mọi công thức \(E\).
  \item Hội: \(E_1\land E_2\land\cdots\land E_n\). Lưu ý rằng phép hội được
  tính từ trái sang phải:
  \[
    E_1\land E_2\land E_3=(E_1\land E_2)\land E_3.
  \]
  \item Tuyển: \(E_1\lor E_2\lor\cdots\lor E_n\). Nhận xét tương tự cũng áp
  dụng cho phép này.
  \item Kéo theo: \(E_1\Rightarrow E_2\). Khác với hai phép toán trước, phép
  này được tính từ phải sang trái:
  \[
    E_1\Rightarrow E_2\Rightarrow E_3
    \quad\text{có nghĩa là}\quad
    E_1\Rightarrow(E_2\Rightarrow E_3).
  \]
  \item Tương đương: \(E_1\equiv E_2\equiv\cdots\equiv E_n\). Biểu thức này
  bằng
  \[
    (E_1\equiv E_2)\land(E_2\equiv E_3)\land\cdots\land(E_{n-1}\equiv E_n).
  \]
\end{itemize}

Các phép toán được liệt kê theo thứ tự từ độ ưu tiên cao nhất đến thấp nhất.

Một công thức \(E\) được gọi là \term{hợp lệ} trong mô hình do \(\mathcal X\)
xác định nếu sau khi thay các phần tử tùy ý của \(H\) cho những biến xuất hiện
trong \(E\), giá trị của công thức luôn là \(1\). Ngược lại, công thức được gọi
là \term{không hợp lệ}.

Nhiệm vụ của bạn là, với một đồ thị \(\mathcal X\) cho trước, xác định những
công thức nào trong một tập công thức đã cho là hợp lệ và những công thức nào
không hợp lệ.

\section*{Dữ liệu vào}

Dòng đầu tiên chứa hai số nguyên \(N\) và \(M\), cách nhau bởi một dấu cách:
số đỉnh \((1\le N\le 100)\) và số cạnh \((0\le M\le 5000)\) của
\(\mathcal X\). \(M\) dòng tiếp theo, mỗi dòng chứa hai số nguyên \(s_i\) và
\(t_i\): đỉnh đầu và đỉnh cuối của cạnh thứ \(i\). Dòng tiếp theo chứa
\(K\) \((1\le K\le 20)\), số công thức cần xử lý. \(K\) dòng sau đó, mỗi dòng
chứa một công thức.

Một công thức được biểu diễn bằng một xâu gồm các token
\[
  \texttt{0},\ \texttt{1},\ \texttt{A},\ldots,\texttt{Z},\ \texttt{(},\
  \texttt{)},\ \texttt{\~{}},\ \texttt{\&},\ \texttt{|},\ \texttt{=>},\
  \texttt{=}.
\]
Năm token cuối lần lượt biểu diễn \(\neg,\land,\lor,\Rightarrow\) và
\(\equiv\). Các token có thể được phân tách bởi số lượng dấu cách tùy ý. Không
dòng nào dài quá \(254\) ký tự. Mọi công thức trong tệp vào đều đúng cú pháp.
Bạn cũng có thể giả sử rằng số \(H=|H|\), tức số phần tử của \(H\), không vượt
quá \(100\), và
\[
  \sum_{1\le j\le K} H^{v_j}\le 10^6,
\]
trong đó \(v_j\) là số biến khác nhau được dùng trong công thức thứ \(j\).

\section*{Dữ liệu ra}

Tệp ra phải chứa \(K\) dòng, mỗi dòng ứng với một công thức. Ghi vào dòng thứ
\(j\) của tệp ra một trong hai từ \texttt{valid} hoặc \texttt{invalid}.

\section*{Ví dụ 1}

\noindent\begin{minipage}[t]{0.48\linewidth}
\textbf{\texttt{logic.in}}
\begin{verbatim}
1 0
6
1=0
X|~X
A=>B=>C = (A&B)=>C
~~X => X
X => ~~X
(X => Y) = (Y | ~X)
\end{verbatim}
\end{minipage}\hfill
\begin{minipage}[t]{0.45\linewidth}
\textbf{\texttt{logic.out}}
\begin{verbatim}
invalid
valid
valid
valid
valid
valid
\end{verbatim}
\end{minipage}

\section*{Ví dụ 2}

\noindent\begin{minipage}[t]{0.48\linewidth}
\textbf{\texttt{logic.in}}
\begin{verbatim}
6 6
1 2
2 3
2 4
3 5
4 5
5 6
11
1=0
X|~X
A=>B=>C = (A&B)=>C
~~X => X
X => ~~X
(X => Y) = (Y | ~X)
A&(B|C) = A&B|A&C
(X=>A)&(Y=>A) => X|Y=>A
X = ~~X
~X=~~~X
~X = (X => 0)
\end{verbatim}
\end{minipage}\hfill
\begin{minipage}[t]{0.45\linewidth}
\textbf{\texttt{logic.out}}
\begin{verbatim}
invalid
invalid
valid
invalid
valid
invalid
valid
valid
invalid
valid
valid
\end{verbatim}
\end{minipage}

\end{document}
